% Detailed proofs for 04qed-one-loop-ward-renormalization

\begin{derivation}[从\eqnrefs{eq:qed-renormalized-electron-inverse-dp53}到\eqnrefs{eq:qed-onshell-selfenergy-conditions-dp53}]
正文把极点问题化成一个真正的一维标量问题。固定未来指向单位类时四矢量 $u^\mu$，令
\begin{equation*}
 p^\mu=\lambda u^\mu,
 \qquad
 P_+(u)=\frac{\id_4+\slashed u}{2}.
\end{equation*}
因为 $u^2=1$，Clifford 关系给
\begin{equation*}
 (\slashed u)^2=\id_4,
 \qquad
 P_+\slashed u=\slashed uP_+=P_+.
\end{equation*}
因此自由逆传播子投影到正能子空间后只有一个标量因子：
\begin{align*}
 P_+(\slashed p-\me c)P_+
 &=P_+(\lambda\slashed u-\me c)P_+\\
 &=(\lambda-\me c)P_+.
\end{align*}

一般 Lorentz 协变、宇称偶的电子自能可写成
\begin{equation*}
 \Sigma_R(p)=A_R(p^2)\slashed p+B_R(p^2)\me c,
\end{equation*}
其中 $A_R,B_R$ 是标量函数。沿 $p=\lambda u$ 投影：
\begin{align*}
 P_+\Sigma_R(\lambda u)P_+
 &=\left[\lambda A_R(\lambda^2)+\me c\,B_R(\lambda^2)\right]P_+.
\end{align*}
这说明正文定义的
\begin{equation*}
 \sigma_R(\lambda)
 =\frac12\Tr\!\left[P_+\Sigma_R(\lambda u)\right]
\end{equation*}
确实就是正能二维子空间上的自能本征值，因为 $\Tr P_+=2$。于是完整逆传播子满足
\begin{equation*}
 P_+S_R^{-1}(\lambda u)P_+
 =d_+(\lambda)P_+,
\end{equation*}
其中
\begin{equation*}
 d_+(\lambda)
 \equiv\lambda-\me c-\sigma_R(\lambda).
\end{equation*}

物理单电子极点位于 $\lambda=\me c$ 的意思是正能投影逆传播子在那里有零本征值：
\begin{equation*}
 d_+(\me c)=0.
\end{equation*}
代入定义即得
\begin{equation*}
 \sigma_R(\me c)=0.
\end{equation*}
这一步只固定零点位置，也就是物理质量；还没有固定极点留数。

再把 $d_+$ 在零点附近作普通一元 Taylor 展开。令
\begin{equation*}
 \delta\lambda\equiv\lambda-\me c,
\end{equation*}
则
\begin{align*}
 d_+(\me c+\delta\lambda)
 &=d_+(\me c)
 +d_+'(\me c)\delta\lambda
 +\frac12d_+''(\me c)(\delta\lambda)^2+\cdots\\
 &=\left[1-\sigma_R'(\me c)\right]\delta\lambda
 +O((\delta\lambda)^2).
\end{align*}
因此正能极点部分的传播子为
\begin{equation*}
 S_R(\lambda u)P_+
 =\frac{\ii Z_{\rm pole}P_+}{\delta\lambda+\ii0^+}
 +S_{R,\mathrm{reg}}(\lambda;u)P_+,
\end{equation*}
其中正则项在 $\delta\lambda=0$ 处有限，而留数由
\begin{equation*}
 Z_{\rm pole}
 =\frac{1}{d_+'(\me c)}
 =\frac{1}{1-\sigma_R'(\me c)}
\end{equation*}
唯一确定。在壳方案把重整化电子场定义成单位极点留数，即
\begin{equation*}
 Z_{\rm pole}=1,
\end{equation*}
于是
\begin{equation*}
 \sigma_R'(\me c)=0.
\end{equation*}
与质量条件合起来正是正文\eqnrefs{eq:qed-onshell-selfenergy-conditions-dp53}。

这个推导同时说明了两条反项条件为什么独立：质量反项调节 $d_+$ 的零点位置，场强反项调节同一零点的一阶导数；它们分别控制 Laurent 展开的极点位置和留数。
\end{derivation}

\begin{derivation}[从\eqnrefs{eq:qed-one-loop-selfenergy-vertex-kernels-dp68}到\eqnrefs{eq:qed-one-loop-ward-kernel-dp53}]

正文\eqnrefs{eq:qed-one-loop-selfenergy-vertex-kernels-dp68} 已把零动量转移处的自能图和顶角图写成同一内部电子传播核的两个积分。关键步骤是证明：对自能图的内部电子传播核关于外动量求导，恰好等价于在该电子线上插入一个 \(\gamma^\mu\) 顶角。

自由约化传播核满足
\begin{equation*}
(\slashed k-\me c)\mathsf S_0(k)=\id_4.
\end{equation*}
对一般矩阵恒等式 \(M M^{-1}=I\) 求导并左乘 \(M^{-1}\)，得到
\begin{equation*}
\partial_xM^{-1}
=-M^{-1}(\partial_xM)M^{-1}.
\end{equation*}
取 \(M(p)=\slashed p-\slashed\ell-\me c\)，则
\begin{equation*}
\frac{\partial\mathsf S_0(p-\ell)}{\partial p_\mu}
=-\mathsf S_0(p-\ell)\gamma^\mu\mathsf S_0(p-\ell).
\end{equation*}
因此对正文自能核逐项求导可得
\begin{align*}
\frac{\partial\Sigma_{\rm 1L}(p)}{\partial p_\mu}
&=-\ii g_{\rm em}^2
\int_{\ell}^{(d)}
\frac{\gamma^\alpha
[\partial\mathsf S_0(p-\ell)/\partial p_\mu]
\gamma_\alpha}
{\ell^2+\ii0^+}\\
&=+\ii g_{\rm em}^2
\int_{\ell}^{(d)}
\frac{\gamma^\alpha\mathsf S_0(p-\ell)
\gamma^\mu\mathsf S_0(p-\ell)\gamma_\alpha}
{\ell^2+\ii0^+}.
\end{align*}
第二行与正文\eqnrefs{eq:qed-one-loop-selfenergy-vertex-kernels-dp68} 中的顶角积分逐因子相同，只差整体负号，所以
\begin{equation*}
\Lambda_{\rm 1L}^{\mu}(p,p)
=-\frac{\partial\Sigma_{\rm 1L}(p)}{\partial p_\mu},
\end{equation*}
即正文\eqnrefs{eq:qed-one-loop-ward-kernel-dp53}。
\end{derivation}

\begin{derivation}[从\eqnrefs{eq:qed-one-loop-ward-kernel-dp53}到\eqnrefs{eq:qed-Z1-Z2-dp53}]

\eqnrefs{eq:qed-one-loop-ward-kernel-dp53} 已经说明一圈电子自能的动量依赖与零动量顶角修正并非独立。把同一关系施加到重整化后的二点函数和顶角后，剩余的局域反项也必须满足相同的系数关系。

把电子二点与电磁顶角的反项分别写成
\begin{align*}
S_R^{-1}(p)
&=\slashed p-\me c-\Sigma_{\rm 1L}(p)
+\delta_2\slashed p-\delta_m c,\\
\Gamma_R^\mu(p,p)
&=\gamma^\mu+\Lambda_{\rm 1L}^\mu(p,p)
+\delta_1\gamma^\mu.
\end{align*}
对第一式关于 \(p_\mu\) 求导，质量反项不贡献：
\begin{equation*}
\frac{\partial S_R^{-1}}{\partial p_\mu}
=\gamma^\mu
-\frac{\partial\Sigma_{\rm 1L}}{\partial p_\mu}
+\delta_2\gamma^\mu.
\end{equation*}
零动量 Ward 恒等式给出
\begin{equation*}
\Gamma_R^\mu(p,p)
=\frac{\partial S_R^{-1}(p)}{\partial p_\mu}.
\end{equation*}
再用正文\eqnrefs{eq:qed-one-loop-ward-kernel-dp53} 比较两边，树级项和一圈图项已经逐项相同，因此只能有
\begin{equation*}
\delta_1\gamma^\mu=\delta_2\gamma^\mu.
\end{equation*}
于是 \(\delta_1=\delta_2\)，等价地
\begin{equation*}
Z_1=Z_2,
\end{equation*}
得到正文\eqnrefs{eq:qed-Z1-Z2-dp53}。
\end{derivation}

\begin{derivation}[从\eqnrefs{eq:qed-bare-renormalized-charge-match-dp68}到\eqnrefs{eq:qed-charge-Z3-dp53}]
关键步骤不是把 \(Z_1=Z_2\) 代进一个已经写好的等式，而是先说明该等式本身从哪里来。裸 QED 的电子--光子相互作用写成
\begin{equation*}
\mathcal L_{\rm int}^{(0)}
=-q_0\,\bar\psi_0\gamma^\mu A_{0\mu}\psi_0.
\end{equation*}
定义重整化场
\begin{equation*}
\psi_0=Z_2^{1/2}\psi_R,
\qquad
\bar\psi_0=Z_2^{1/2}\bar\psi_R,
\qquad
A_{0\mu}=Z_3^{1/2}A_{R\mu},
\end{equation*}
则同一个裸相互作用在重整化场变量中成为
\begin{align*}
\mathcal L_{\rm int}^{(0)}
&=-q_0
\bigl(Z_2^{1/2}\bar\psi_R\bigr)
\gamma^\mu
\bigl(Z_3^{1/2}A_{R\mu}\bigr)
\bigl(Z_2^{1/2}\psi_R\bigr)\\
&=-q_0Z_2Z_3^{1/2}
\bar\psi_R\gamma^\mu A_{R\mu}\psi_R.
\end{align*}
另一方面，重整化拉氏量把同一个局域 Dirac 顶角写成“重整化耦合 + 顶角反项”，即
\begin{equation*}
\mathcal L_{\rm int}^{R+\rm ct}
=-q_R Z_1\,
\bar\psi_R\gamma^\mu A_{R\mu}\psi_R.
\end{equation*}
两式描述的是同一个裸理论，只是变量不同；因为
\(\bar\psi_R\gamma^\mu A_{R\mu}\psi_R\)
是同一个局域算符，其系数必须逐项相等，所以
\begin{equation*}
q_0Z_2Z_3^{1/2}=q_RZ_1,
\end{equation*}
这才得到正文\eqnrefs{eq:qed-bare-renormalized-charge-match-dp68}。

上一推导由 Ward--Takahashi 恒等式得到 \(Z_1=Z_2\)。因此
\begin{equation*}
q_0Z_2Z_3^{1/2}=q_RZ_2.
\end{equation*}
场重整化必须是可逆变量重定义，所以 \(Z_2\neq0\)，约去以后得到
\begin{equation*}
q_0Z_3^{1/2}=q_R,
\qquad
q_0=q_RZ_3^{-1/2},
\end{equation*}
即正文\eqnrefs{eq:qed-charge-Z3-dp53}。

这一步的结构意义现在是可见的：电子外腿的两个 \(Z_2^{1/2}\) 本来确实进入裸顶角，但局域 \(U(1)\) 对称性又强迫顶角反项满足 \(Z_1=Z_2\)，两者严格抵消；因而独立的电荷紫外重整化只能由光子二点函数中的 \(Z_3\) 承担。
\end{derivation}
